\section{Match Transfer Claims to the Generator}\label{sec:evaluation}\label{sec:threats}
\subsection{Different tests rule out different explanations}
A fingerprint is generated by a model interacting with tasks, prompts, tools, browser state and serving infrastructure. A holdout establishes portability along the axis it changes: new tasks, browser configurations, network topologies and hardware are distinct deployment tests. Supplement S2 maps each dependency to its supporting holdout.

UI Traces fixes a browser/action harness while testing tasks and environments. FP-Agent combines behavior with static browser attributes; shared visitor proxies are not authenticated independent people or hosts \cite{uitraces2026,wang2026fpagent}. Task separation can remove memorization of those task instances while retaining configuration identifiers. This is not proof of an invalid original split; it bounds portability. The traffic benchmark separately tests generating topology, so its shift is a different test rather than a stronger score on the same estimand \cite{pouryousef2026traffic}.

LeakyLMs illustrates the boundary for system-property rather than agent identity. Its deployment probe observes streaming events through an ordinary client API. Its architecture branch additionally assumes known GPU type, same-class calibration and single-GPU execution \cite{majidi2026leakylms}. Remote architecture inference therefore does not establish hardware-agnostic inference. The joint search in its remote example ranks the true configuration 28/2068, distinct from a conditional parameter search or a near match; those outcomes must not be collapsed into universal exact recovery. Source-reported local architecture generalization and remote deployment signatures remain positive evidence under their own assumptions.

\subsection{Representation increment is not portability}
Within a fixed observer and test population, coarse counts versus sequence features test whether the latter improve a specified predictor. WRA already uses ordered domains before adding timing; Kang changes class sets across some conditions \cite{jeong2026wra,kang2026whoseagent}. These test different contrasts from an order-only effect. The supplementary archive example fixes labels and units and finds a 0.3447-nat sequence increment under archive-receiver ordering. Its small common-task control leaves task-distribution explanations unresolved (Supplement S4).

Adjacent coding-agent PR attribution depends on repositories, submission pathways and artifact generation. Its mixed-feature result does not isolate an auxiliary-only C3 view; inspected code computes weighted F1 and reduces features before cross-validation, but the exact paper-run binding remains unverified \cite{ghaleb2026codingfingerprint}. We neither rewrite its headline score nor quantify optimism. WildAuth is adjacent under the same rule; input/partition ambiguities do not supply a model-independent response-only privacy claim \cite{wang2026dialogue}.

\subsection{Resulting claim hierarchy}
Task-specific identification, task-disjoint identification in a fixed harness, configuration/topology transfer and hardware-calibrated architecture inference require different evidence. Move between them by changing the relevant generator while preserving target labels and the receiver contract. For a representation increment, compare representations on identical units and disclose predictor/selection restrictions. This claim-to-holdout rule lets a reader reuse each author test for the shift it measures.
